\documentclass[11pt,a4paper]{article}
\usepackage[T1]{fontenc}
\usepackage[utf8]{inputenc}
\usepackage{lmodern,amsmath,amssymb,amsthm,mathtools,booktabs,microtype}
\usepackage[margin=27mm]{geometry}
\usepackage{xcolor}
\definecolor{ink}{HTML}{392E46}
\usepackage[colorlinks=true,linkcolor=ink,citecolor=ink,urlcolor=ink]{hyperref}
\usepackage{enumitem}
\setlist{itemsep=2pt,topsep=4pt}
\newtheorem{theorem}{Theorem}[section]
\newtheorem{lemma}[theorem]{Lemma}
\newtheorem{corollary}[theorem]{Corollary}
\newtheorem{proposition}[theorem]{Proposition}
\theoremstyle{remark}
\newtheorem{remark}[theorem]{Remark}
\newcommand{\U}{\mathcal U}
\newcommand{\J}{\mathcal J}
\newcommand{\F}{\mathcal F}
\newcommand{\D}{\mathcal D}
\newcommand{\wt}{\boldsymbol w}
\newcommand{\ind}{\mathbf 1}
\newcommand{\bin}[2]{\binom{#1}{#2}}
\title{Ultra log-concavity for increasing families\protect\\ generated by singletons and pairs}
\author{Anonymous}
\date{4 November 2026\\[2pt] \small Version 0.1.0-candidate. DOI: \href{https://doi.org/10.5281/zenodo.23935650}{10.5281/zenodo.23935650}}
\begin{document}
\maketitle
\begin{abstract}
Pick a graph and count, for each $k$, the $k$-vertex sets that contain at least one edge. Xie and Xu conjectured that these counts, divided by $\binom nk$, form a log-concave sequence: the strong ``ultra log-concave'' form of an older theorem of Horrocks. We prove this conjecture (Xie--Xu Conjecture 4.1) in a weighted form. Let $\U$ be an increasing family of subsets of an $n$-element set whose minimal members have size at most two. For arbitrary nonnegative element weights, the weighted rank sequence of $\U$, divided by $\binom nk$, is log-concave. The proof first establishes a stronger two-colouring statement: the binomially normalised number of $k$-sets $S$ such that both $S$ and its complement contain a witness is symmetric and unimodal, and this survives fixing any coordinates. It combines a marked-family consequence of Katona's intersection-shadow theorem, a critical-element double count, and Kahn and Neiman's antipodal-pairs criterion. Applying the graph theorem, we prove a weighted strengthening of Xie and Xu's Conjecture 4.4 for additive $2$-matroids of rank at most four or on at most seven elements; the unrestricted conjecture remains open. The witness-size bound is sharp: an unweighted five-element family with mixed singleton and triple witnesses fails ULC. Even two disjoint triple witnesses fail for suitable positive weights, and six is the smallest ground-set size admitting a positively weighted counterexample whose minimal witnesses all have size three. A variance bound for conditioned Bernoulli occupancy follows.
\end{abstract}

\section{Statement and position in the literature}

\paragraph{Reader's guide.} A nonnegative sequence $a_0,\dots,a_n$ is \emph{log-concave} if $a_k^2\ge a_{k-1}a_{k+1}$ for $1\le k\le n-1$. If it has no internal zeros, it is then unimodal (it rises and falls with no second hump), and its successive ratios $a_{k+1}/a_k$ are nonincreasing wherever both terms are positive. It is \emph{ultra log-concave of order $n$} (ULC) if $(a_k/\binom nk)_{k=0}^n$ is log-concave and has no internal zeros. Binomial probability sequences attain equality. Informally, ULC says a distribution on $\{0,\dots,n\}$ is at least as concentrated as a binomial one; for example, a ULC probability distribution with mean $m$ has variance at most $m(1-m/n)$ (Corollary~\ref{cor:variance}). ULC is the property in Mason's famous conjecture on matroid independent sets, now a theorem (Section~\ref{sec:related}). Here it is proved for counts of \emph{dependent} vertex sets in graphs. The paper proceeds as follows.
\begin{enumerate}
\item Section~1 states the main theorem (Theorem~\ref{thm:main}) and the stronger two-colouring theorem (Theorem~\ref{thm:apu}) from which it follows.
\item Section~2 proves two counting lemmas about downsets. The new one, Lemma~\ref{lem:marked}, is a short consequence of Katona's shadow theorem.
\item Section~3 proves Theorem~\ref{thm:apu} by counting ``critical'' elements.
\item Section~4 converts Theorem~\ref{thm:apu} into Theorem~\ref{thm:main} through the Kahn--Neiman criterion, quoted exactly.
\item Section~5 gives consequences, counterexamples showing the hypotheses are needed, and limits. Section~6 treats the $2$-matroid extension. Section~7 records verification and status.
\end{enumerate}

\paragraph{A worked example.} Take two disjoint edges $ab$ and $cd$ on $n=4$ vertices. The dependent sets are those containing $\{a,b\}$ or $\{c,d\}$, so the counts are $(d_0,\dots,d_4)=(0,0,2,4,1)$. Dividing by $\binom4k$ gives $(0,0,\tfrac13,1,1)$. At $k=3$ the ULC inequality reads $1^2\ge\tfrac13\cdot1$, and at $k=2$ it reads $(\tfrac13)^2\ge0$; both hold. The two-colouring count $b_2$ in Theorem~\ref{thm:apu} counts the two sets $\{a,b\},\{c,d\}$, whose complements also contain an edge.

\paragraph{Key terms.} \emph{Witness}: a single element (a loop) or a pair (an edge) whose presence makes a set dependent. \emph{Increasing family}: a family closed under taking supersets; here, the dependent sets. \emph{Doubly dependent set}: a set that contains a witness, as does its complement. \emph{Critical element}: an element whose removal makes a dependent set independent. \emph{APP/CAPP}: the Kahn--Neiman antipodal-pairs property and its conditional form (Section~4). \emph{CAPU}: the stronger conditional antipodal-pairs unimodality proved here.

Let $V$ be a finite set of size $n$. A family $\U\subseteq2^V$ is \emph{increasing} if $S\in\U$ and $S\subseteq T$ imply $T\in\U$. We say that $\U$ is \emph{generated by singletons and pairs} if there is a collection
\[
\mathcal H\subseteq\binom V1\cup\binom V2
\quad\text{such that}\quad
\U=\{S\subseteq V:\text{some }H\in\mathcal H\text{ is contained in }S\}.
\]
Members of $\mathcal H$ will be called witnesses. Redundant witnesses may be removed. Equivalently, $\U$ is the dependent-set family of a graph in which loops are permitted: a loop is a singleton witness. The empty family and $2^V$ will be treated separately when needed.

For $w_v\geq0$, put $w(S)=\prod_{v\in S}w_v$, with $w(\varnothing)=1$, and define
\[
A_k(\U;\wt)=\sum_{\substack{S\in\U\\|S|=k}}w(S),
\qquad
D_{\U}(z;\wt)=\sum_{k=0}^n A_k(\U;\wt)z^k.
\]
Our normalisation always uses the ambient cardinality $n$, including any isolated or zero-weight elements.

\begin{theorem}[Weighted rank inequality]\label{thm:main}
If $\U$ is generated by singletons and pairs, then for every choice of nonnegative element weights and every $1\leq k\leq n-1$,
\begin{equation}\label{eq:main}
k(n-k)A_k^2\ \geq\ (k+1)(n-k+1)A_{k-1}A_{k+1}.
\end{equation}
Thus $(A_k/\binom nk)_{k=0}^n$ is log-concave. The same inequalities hold for $\U=2^V$ and $\U=\varnothing$.
\end{theorem}

For an ordinary finite simple graph $G$, let $d_k(G)$ count its $k$-vertex subsets containing an edge. Taking $w_v=1$ gives
\begin{equation}\label{eq:graph}
\left(\frac{d_k(G)}{\binom nk}\right)^2\geq
\frac{d_{k-1}(G)}{\binom n{k-1}}
\frac{d_{k+1}(G)}{\binom n{k+1}}.
\end{equation}
This is precisely Conjecture 4.1 in the 10 October 2026 version of Xie and Xu \cite{XX}. Their dependence polynomial is $(1+z)^n-I(G,z)$, where $I$ is the independence polynomial. Horrocks \cite{H02} proved ordinary log-concavity. Xie and Xu, citing that result, posed the binomially normalised inequality as open and proved it for two classes of graphs: $(K_2\cup2K_1)$-free graphs and graphs with an independent set of size $n-2$. Our argument imposes no class restriction. We credit Horrocks with ordinary log-concavity only. His paper's abstract states that result, but we could not obtain its full text, so we have not compared proof techniques.

We also strengthen a theorem of Horrocks \cite{H00} about \emph{doubly dependent sets}. Define
\[
\D=\{S\subseteq V:S\in\U,\ V\setminus S\in\U\},
\qquad b_k=|\D\cap\binom Vk|.
\]
\begin{theorem}[Normalised two-colouring counts]\label{thm:apu}
For $\U$ as above, the sequence $(b_k/\binom nk)_{k=0}^n$ is symmetric and unimodal, with its maximum at the middle rank or ranks. In particular,
\begin{equation}\label{eq:apu}
k b_k\ \geq\ (n-k+1)b_{k-1}
\qquad (1\leq k\leq\lfloor n/2\rfloor).
\end{equation}
\end{theorem}

Horrocks's unnormalised unimodality theorem does not supply the factor in \eqref{eq:apu}. That factor permits the use of the antipodal-pairs criterion of Kahn and Neiman \cite{KN}. We state every external theorem used in the proof below. The key step is the marked-family comparison in Lemma~\ref{lem:marked} and its use in the critical-element count. We do not claim that the constituent shadow theorem or the Kahn--Neiman implication is new.

\subsection{Related work}\label{sec:related}
Mason's ultra log-concavity conjecture for matroid independent sets was proved by Anari, Liu, Oveis Gharan and Vinzant \cite{ALOV} using completely log-concave polynomials, and by Br\"and\'en and Huh \cite{BH} using Lorentzian polynomials. A homogeneous Lorentzian polynomial must have M-convex support. The natural homogenisation
\[
F_{\U}(x_0,\boldsymbol x)=\sum_{S\in\U}x_0^{n-|S|}\prod_{v\in S}x_v
\]
need not satisfy this requirement. For two disjoint edges $ab,cd$, its support contains the monomials $x_0^2x_ax_b$ and $x_0^2x_cx_d$, but no exchange of $a$ for $c$ or $d$ produces a supported monomial. Thus this natural homogenisation is not Lorentzian, and we instead use the antipodal-pairs criterion of Kahn and Neiman \cite{KN}. In topological language, the downset of complements of dependent sets is the \emph{noncover complex} of the graph \cite{KK}. The dependence polynomial is $(1+z)^n-I(G,z)$. Ultra log-concave sequences are closed under convolution \cite{Liggett}; adding an isolated element convolves the rank sequence with $(1,w_v)$, which is why the normalisation order matters (Remark~\ref{rem:order}).

\section{Two shadow comparisons}

A family is a \emph{downset} if it contains every subset of each of its members. For an $r$-uniform family $\F$, its $\ell$-shadow is
\[
\partial_\ell\F=\{L:|L|=\ell,\ L\subseteq F\text{ for some }F\in\F\}.
\]

\begin{lemma}[Elementary rank comparison]\label{lem:ordinary}
Let $\F$ be a downset on an $N$-element ground set. If $0\leq\ell\leq r$ and $\ell+r\geq N$, then $|\F_\ell|\geq|\F_r|$, where $\F_j$ is the $j$th rank.
\end{lemma}
\begin{proof}
The result is immediate if $r>N$. Otherwise, counting pairs $(L,R)$ with $L\subseteq R\in\F_r$ and $|L|=\ell$ gives
\[
\binom r\ell|\F_r|\leq\binom{N-\ell}{r-\ell}|\F_\ell|.
\]
Consequently $|\F_\ell|\geq\binom N\ell\binom Nr^{-1}|\F_r|$. The binomial ratio is at least one because $\ell\leq r$ and $\ell+r\geq N$.
\end{proof}

We use the following form of Katona's intersection-shadow theorem \cite{Katona}, stated on the first page of Frankl's author manuscript \cite{Frankl}, which also provides a proof. If an $r$-uniform family $\F$ is $t$-intersecting, meaning $|F\cap F'|\geq t$ for all $F,F'\in\F$, then
\begin{equation}\label{eq:katona}
|\partial_{r-t}\F|\geq|\F| \qquad (1\leq t\leq r).
\end{equation}
The endpoint $t=r$ is immediate since a nonempty such family has one member. We use no equality classification.

\begin{lemma}[Marked downset comparison]\label{lem:marked}
Let $\J$ be a downset on a set $W$ of size $m$, and suppose
\begin{equation}\label{eq:nounion}
I\cup J\ne W\qquad\text{for every }I,J\in\J.
\end{equation}
Fix a set $A\subseteq W$. If $0\leq a\leq b$ and $a+b=m$, then
\begin{equation}\label{eq:marked}
|\{I\in\J_a:I\cap A\ne\varnothing\}|
\ \geq\
|\{I\in\J_b:I\cap A\ne\varnothing\}|.
\end{equation}
\end{lemma}
\begin{proof}
If $A=\varnothing$ or $a=b$, the assertion is immediate. If $a=0$, both sides vanish: the left side cannot hit $A$, and \eqref{eq:nounion} excludes $W$ from $\J$. Hence assume $1\leq a<b$ and order $A=\{u_1,\ldots,u_q\}$.

Partition the sets meeting $A$ according to their first element $u_i$ in this order. After deleting $u_i$, the $i$th class is a downset
\[
\F_i=\{T\subseteq W\setminus\{u_1,\ldots,u_i\}:T\cup\{u_i\}\in\J\}.
\]
Its contributions to the two sides are $|\F_{i,a-1}|$ and $|\F_{i,b-1}|$. For $i\geq2$, its ground set has $m-i\leq m-2$ elements, and
\[
(a-1)+(b-1)=m-2.
\]
Lemma~\ref{lem:ordinary} therefore gives the required comparison for every $i\geq2$.

For $i=1$, the ground set of $\F_1$ is $W\setminus\{u_1\}$, of size $m-1$. Let $\mathcal B=\F_{1,b-1}$, a $(b-1)$-uniform family. If $T,T'\in\mathcal B$, then \eqref{eq:nounion}, applied to $T\cup\{u_1\}$ and $T'\cup\{u_1\}$, implies $|T\cup T'|\leq m-2$. Thus
\[
|T\cap T'|\geq2(b-1)-(m-2)=b-a.
\]
Apply \eqref{eq:katona} with $r=b-1$ and $t=b-a$; here $t\ge1$ because $a<b$, and $t\le r$ because $a\ge1$. Its shadow has rank $r-t=a-1$, so
\[
|\F_{1,a-1}|\geq|\partial_{a-1}\mathcal B|
\geq|\mathcal B|=|\F_{1,b-1}|.
\]
Summing over the disjoint first-element classes proves \eqref{eq:marked}.
\end{proof}

\section{Critical elements and normalised unimodality}

For $S\in\U$, call $v\in S$ \emph{critical} if $S\setminus\{v\}\notin\U$. Put
\[
c(S)=|\{v\in S:S\setminus\{v\}\notin\U\}|,
\qquad C_j=\sum_{S\in\D\cap\binom Vj}c(S).
\]
No bound on $c(S)$ will be needed.

\begin{lemma}[Boundary identity]\label{lem:boundary}
For any increasing family $\U$ and $1\leq k\leq n$,
\begin{equation}\label{eq:boundary}
k b_k-(n-k+1)b_{k-1}=C_k-C_{n-k+1}.
\end{equation}
\end{lemma}
\begin{proof}
Count inclusions $S\subset T$ with $S,T\in\D$, $|S|=k-1$ and $|T|=k$. A $T\in\D_k$ has $k-c(T)$ admissible deletions: its complement remains dependent after adding a vertex. This gives $k b_k-C_k$.

For an $S\in\D_{k-1}$, adding a vertex preserves dependence of $S$. Exactly $c(V\setminus S)$ additions destroy dependence of the complement. The count is therefore $(n-k+1)b_{k-1}-C_{n-k+1}$, where complementation reindexes the sum. Equating the counts proves the identity.
\end{proof}

\begin{proof}[Proof of Theorem~\ref{thm:apu}]
Symmetry follows from complementation. If $\U$ is empty or full, the conclusion is immediate. Otherwise $\varnothing\notin\U$, so $b_0=0$ and \eqref{eq:apu} is automatic for $k=1$. Fix $2\leq k\leq\lfloor n/2\rfloor$. By Lemma~\ref{lem:boundary}, it suffices to prove $C_k\geq C_{n-k+1}$.

For each $v\in V$, let $W=V\setminus\{v\}$ and define
\[
\J_v=\{I\subseteq W:I\notin\U,\ W\setminus I\in\U\}.
\]
This is a downset. It also satisfies \eqref{eq:nounion}: if $I,J\in\J_v$ and $I\cup J=W$, then $W\setminus I\subseteq J$, so the increasing property would give $J\in\U$, a contradiction.

The contribution of $v$ to $C_j$ is the number of $I\in\J_v$ of size $j-1$ for which $I\cup\{v\}\in\U$. If $\{v\}$ is a witness, this last condition is automatic. Lemma~\ref{lem:ordinary} compares the ranks $k-1$ and $n-k$ of $\J_v$, since their sum is $|W|=n-1$.

If $\{v\}$ is not a witness, put
\[
N_v=\{u\in W:\{u,v\}\in\mathcal H\}.
\]
Because $I\notin\U$, the condition $I\cup\{v\}\in\U$ is equivalent to $I\cap N_v\ne\varnothing$. Lemma~\ref{lem:marked}, with $a=k-1$ and $b=n-k$, compares the two contributions. Thus each vertex contributes at least as much to $C_k$ as to $C_{n-k+1}$. Summing proves \eqref{eq:apu}.
\end{proof}

\section{Conditional antipodal pairs and weighted ultra log-concavity}

For a probability measure $\mu$ on $2^V$, write
\[
\alpha_j(\mu)=\binom nj^{-1}\sum_{|S|=j}\mu(S)\mu(V\setminus S).
\]
The measure has \emph{antipodal-pairs unimodality} if this sequence is unimodal. It has the conditional version, abbreviated CAPU, if every measure obtained by a positive-probability fixing of any subset of the coordinates has this property on its remaining coordinates. Symmetry of $\alpha_j$ is automatic.

Kahn and Neiman define, for $\mu$ on $2^V$ with $|V|=2k$, the \emph{antipodal pairs property} (APP) as $\alpha_k(\mu)\ge\alpha_{k-1}(\mu)$. A measure has the \emph{conditional antipodal pairs property} (CAPP) if every measure obtained from it by conditioning on the values of some $n-2k$ variables (for some $k$) has the APP. Their main result \cite[Theorem 1]{KN} reads: ``For measures without internal zeros in their rank sequence, the CAPP implies ULC.'' They also note \cite[Corollary 12]{KN} that every CAPU measure is ULC, since CAPU trivially implies CAPP. We use Theorem~1 in this form. Its hypotheses are checked in the proof of Theorem~\ref{thm:main}: (i) CAPP, which follows from the CAPU property of Theorem~\ref{thm:capu}; (ii) no internal zeros in the \emph{rank sequence}. Positivity of $\mu$ on every set is not required. Throughout, ``CAPU'' denotes the stronger property we prove and ``CAPP'' the property the criterion needs.

\begin{theorem}[Conditional weighted form]\label{thm:capu}
For positive weights, the probability measure
\begin{equation}\label{eq:measure}
\mu(S)=\frac{w(S)\ind_{\{S\in\U\}}}{Z},\qquad
Z=\sum_{T\in\U}w(T),
\end{equation}
is CAPU whenever $\U$ is nonempty and is generated by singletons and pairs. The same holds for $\U=2^V$.
\end{theorem}
\begin{proof}
First consider the original measure. Since $w(S)w(V\setminus S)=w(V)$,
\[
\alpha_j(\mu)=\frac{w(V)}{Z^2}\frac{b_j}{\binom nj}.
\]
Theorem~\ref{thm:apu} proves its unimodality.

Now fix a set $P$ of coordinates to be present and a disjoint set $Q$ to be absent. Write $R=V\setminus(P\cup Q)$. If $P$ already contains a witness, the conditional family is $2^R$. Otherwise a surviving witness either remains a singleton or pair in $R$, or is a pair with one endpoint in $P$ and becomes a singleton in $R$. Witnesses touching $Q$ disappear. Therefore every positive-probability conditional family is again generated by singletons and pairs or is full. Its measure has the form \eqref{eq:measure} using the remaining weights, since the factor $w(P)$ cancels. Apply the first part on $R$.
\end{proof}

\begin{proof}[Proof of Theorem~\ref{thm:main}]
Assume first that all weights are positive and $\U$ is nonempty. Its rank support is an interval: if its smallest member has size $r$, every size from $r$ to $n$ occurs. Theorem~\ref{thm:capu} gives CAPU, hence CAPP, and the rank sequence has no internal zeros. Kahn--Neiman Theorem 1 therefore applies. The rank probabilities are $A_k/Z$, and the common factor cancels, giving \eqref{eq:main}.

For nonnegative weights, replace each $w_v$ by $w_v+\varepsilon$ and let $\varepsilon\downarrow0$. Each side of \eqref{eq:main} is continuous in the weights. The empty family gives zero on both sides. The full family follows by the same CAPU argument (or by Newton's inequalities).
\end{proof}

\begin{remark}[A sharper order]\label{rem:order}
Let $V^+=\{v:w_v>0\}$ and $n^+=|V^+|$. Sets containing a zero-weight element contribute nothing, so $A_k(\U;\wt)=A_k(\U\cap2^{V^+};\wt|_{V^+})$. The family $\U\cap2^{V^+}$ is generated by the witnesses inside $V^+$. Theorem~\ref{thm:main} on $V^+$ therefore gives \eqref{eq:main} with $n$ replaced by $n^+$. This is stronger, because ULC of order $m$ implies ULC of every order $n\ge m$ (cf.\ \cite{Liggett}): the inequality $a_k^2\ge(1+\frac1k)(1+\frac1{n-k})a_{k-1}a_{k+1}$ weakens as $n$ grows. Isolated positive-weight elements cannot be removed in the same way, since they multiply the polynomial by $1+w_vz$. In the unweighted graph case, $n^+=n$.
\end{remark}

\section{Consequences and exact limits}

\subsection{The noncover complex}

For a graph with at least one edge, the family
\[
\operatorname{NC}(G)=\{S\subseteq V:V\setminus S\text{ contains an edge}\}
\]
is a downset with facets $V\setminus e$ as $e$ runs through the edges. Its unweighted rank sequence is the reversal of $(d_k)$. Thus Theorem~\ref{thm:main} proves ultra log-concavity for the face counts of every pure simplicial complex on $n$ labelled ground elements whose facets have size $n-2$. Here labelled ground elements that occur in no facet are allowed; deleting such elements changes the order of binomial normalisation. With singleton witnesses, the facets may have sizes $n-1$ and $n-2$.

The weighted statement also transfers. For positive weights,
\[
\sum_{\substack{S\in\operatorname{NC}(G)\\|S|=k}}w(S)
=w(V)A_{n-k}(\U;\wt^{-1}),
\]
and zero weights follow by continuity. This is an equivalent formulation, not a separate priority claim.

\subsection{Conditioned Bernoulli occupancy}

Assume $n\geq1$. Let $X_v$ be independent Bernoulli variables with probabilities $p_v\in(0,1)$, and suppose $\mathbb P(S\in\U)>0$ for their selected set $S=\{v:X_v=1\}$. Condition on $S\in\U$. The resulting law is \eqref{eq:measure} with $w_v=p_v/(1-p_v)$. Hence its cardinality $K=|S|$ has a binomially normalised log-concave distribution.

\begin{corollary}[Variance bound]\label{cor:variance}
With $m=\mathbb E K$,
\[
\operatorname{Var}(K)\leq m\left(1-\frac mn\right).
\]
\end{corollary}
\begin{proof}
The cases $m=0,n$ are immediate. Let $(p_k)$ denote the law of $K$. Ultra log-concavity says that
\[
\frac{(k+1)p_{k+1}}{(n-k)p_k}
\]
is nonincreasing wherever defined. Consider the distributions
\[
q_k=\frac{(k+1)p_{k+1}}{m},\qquad
r_k=\frac{(n-k)p_k}{n-m}\quad(0\leq k\leq n-1).
\]
Their likelihood ratio $q_k/r_k$ is nonincreasing, with an initial infinite ratio permitted where only $q_k$ is positive. The interval support makes this ordering valid at the endpoints. Therefore $q$ is stochastically dominated by $r$: the difference $q_k-r_k$ changes sign at most once, from positive to negative, and has total sum zero. Comparing means gives
\[
\frac{\mathbb E K^2}{m}-1
\leq\frac{nm-\mathbb E K^2}{n-m}.
\]
Rearrangement yields $n\mathbb E K^2\leq nm+(n-1)m^2$, which is the claimed bound. This is a consequence of ultra log-concavity, not a new general variance principle.
\end{proof}

\subsection{Witness size three already permits failure}

The bound on witness size cannot be increased from two to three in a theorem covering every increasing family. Let $V=\{0,1,2,3,4\}$ and let $\U$ consist of sets that contain $0$ or contain at least three of the other four elements. Its minimal witnesses are $\{0\}$ and the four triples in $\{1,2,3,4\}$. Direct counting gives
\[
(A_0,A_1,A_2,A_3,A_4,A_5)=(0,1,4,10,5,1).
\]
At $k=2$, the two sides of \eqref{eq:main} are respectively
\[
2\cdot3\cdot4^2=96
\quad\text{and}\quad
3\cdot4\cdot1\cdot10=120.
\]
Equivalently $(A_k/\binom5k)=(0,1/5,2/5,1,1,1)$ fails log-concavity at $k=2$. This example has mixed singleton and triple witnesses; it makes no claim about the class of uniform triple systems.

A four-element weighted example is also available. Use witnesses $\{0\}$ and $\{1,2,3\}$ with weights $(1,2,2,2)$. Its coefficients are $(0,1,6,20,8)$, and \eqref{eq:main} fails at $k=2$ because $144<180$.

Even requiring every witness to have size exactly three does not rescue the weighted theorem. On six elements take the two disjoint witnesses $\{1,2,3\}$ and $\{4,5,6\}$, with weights $(1,M,M,3,3,3)$. The nonzero coefficients are
\[
\begin{aligned}
A_3&=M^2+27,& A_4&=9M^2+54M+27,\\
A_5&=54M^2+54M,& A_6&=27M^2.
\end{aligned}
\]
The $k=4$ defect is
\begin{equation}\label{eq:tripledefect}
8A_4^2-15A_3A_5
=162(-M^4+43M^3+33M^2+9M+36).
\end{equation}
For positive integer $M$, this is negative exactly when $M\geq44$. For $M\leq43$ every term in $M^3(43-M)+33M^2+9M+36$ is nonnegative. For $M\geq44$, the negative term has magnitude at least $44M^2$, greater than $33M^2+9M+36$. Thus a matching of two triples already supplies weighted counterexamples.

Six is the smallest ground-set size for such a positively weighted, uniform-three-witness counterexample. With at most four elements there are at most two consecutive nonzero coefficients, so the inequalities are automatic. With five elements, reverse the coefficients and divide by $W=\prod_vw_v$. The result is
\[
1+s_1x+s_2x^2,
\]
where $s_2$ is the weighted edge sum of the graph whose edges are complements of the triple witnesses, using reciprocal weights, and $s_1$ is the sum of reciprocal weights over that graph's nonisolated vertices. If there are $p\leq5$ such vertices, then
\[
s_2\leq\sum_{i<j}x_ix_j
\leq\frac{p-1}{2p}s_1^2\leq\frac25s_1^2,
\]
where the middle step is Maclaurin's inequality for the $p$ reciprocal weights $x_i$ (the edge sum is at most the full second elementary symmetric sum, all terms being positive). Thus
\[
s_2\le\tfrac25s_1^2.
\]
This is precisely the sole nontrivial order-five ultra log-concavity condition, $s_1^2\geq\frac52s_2$. The zero-witness case is trivial.

\subsection{What the theorem does not settle}

Xie and Xu's Conjecture 4.4 concerns all independence systems satisfying their relaxed augmentation axiom, called $2$-matroids. The graph theorem covers the noncover systems described above. Section~\ref{sec:two-matroid} extends the conclusion to all $2$-matroids of rank at most four or on at most seven elements, but does not settle the unrestricted conjecture. The critical-element argument uses the fact that adding an element can create a witness only through a single marked neighbour or a singleton witness. A general relaxed-augmentation system need not have that structure.

CAPU and ultra log-concavity also do not establish pairwise negative correlation, real-rootedness of the dependence polynomial, or a general sampling algorithm. Those properties are not claimed here. For example, in the uniform measure on dependent subsets of two disjoint edges, the two endpoints of an edge are positively correlated: $4/7>(5/7)^2$.

\section{Additive 2-matroids of small rank or ground-set size}
\label{sec:two-matroid}

Following the additive augmentation convention of Xie and Xu \cite{XX}, a
\emph{2-matroid} is a nonempty downset $\mathcal I\subseteq2^V$ such that
\begin{equation}\label{eq:two-augmentation}
S,T\in\mathcal I,\quad |T|\ge |S|+2
\quad\Longrightarrow\quad
S\cup\{x\}\in\mathcal I\text{ for some }x\in T\setminus S.
\end{equation}
Its rank is $\max\{|S|:S\in\mathcal I\}$. For element weights $w_v\ge0$, write
\[
A_j=\sum_{\substack{S\in\mathcal I\\|S|=j}}w(S).
\]

\begin{theorem}\label{thm:two-matroid-small}
The weighted rank sequence of a 2-matroid on $n$ ground elements is ultra
log-concave of ambient order $n$ if its rank is at most four or if $n\le7$.
Thus, in either case,
\[
k(n-k)A_k^2\ge(k+1)(n-k+1)A_{k-1}A_{k+1}
\qquad(1\le k\le n-1).
\]
\end{theorem}

We establish the required antipodal-pairs inequalities analytically. On an
even ground set of size $2s$, put
\[
F_j=\{S\in\mathcal I:|S|=j,\ V\setminus S\in\mathcal I\}.
\]
The inequality to prove is $s|F_s|\ge(s+1)|F_{s-1}|$. On two elements this is
automatic for any downset: if $F_0$ is nonempty, the full two-set is feasible,
so both singletons are feasible and $|F_1|=2$.

\begin{lemma}[Four elements]\label{lem:two-matroid-four}
Every 2-matroid on four elements satisfies $2|F_2|\ge3|F_1|$.
\end{lemma}
\begin{proof}
Assume $F_1\ne\varnothing$, as otherwise the assertion is immediate. A
feasible singleton and its feasible complementary triple show that all four
singletons are feasible. Thus $|F_1|=t$, the number of feasible triples.
View the feasible pairs as a graph.

If $t=1$, augment the singleton outside the feasible triple towards that
triple. This gives a feasible pair joining the outside vertex to a vertex
of the triple; its complementary pair is feasible by downward closure.
Hence $|F_2|\ge2$. If $t=2$, the pair shadows of the two triples contain
all edges of $K_4$ except at most one. At least two perfect matchings remain,
so $|F_2|\ge4$. If $t\ge3$, all six pairs are feasible, so $|F_2|=6$.
These three cases give $2|F_2|\ge3t$.
\end{proof}

\begin{lemma}[Six elements]\label{lem:two-matroid-six}
Every 2-matroid on six elements satisfies $3|F_3|\ge4|F_2|$.
\end{lemma}
\begin{proof}
Let $R$ be the simple graph with edge set $F_2$, and put $r=|E(R)|=|F_2|$.
Let $m$ count unordered feasible partitions of $V$ into two triples. Thus
$|F_3|=2m$, and the target is $3m\ge2r$. The case $r=0$ is immediate.

Let $g$ count unordered balanced partitions with an $R$-edge in each part.
Every such partition is feasible: if $e\subseteq A$ and $f\subseteq B$ are
edges, then $A\subseteq V\setminus f$ and $B\subseteq V\setminus e$, and
the complementary four-set of each $R$-edge is feasible. Hence $m\ge g$.

Suppose first that no edge of $R$ is a vertex cover. Then every edge has an
edge in its complement. For the doubly dependent counts of the ordinary
graph $R$, this says $b_2(R)=r$ and $b_3(R)=2g$. Theorem~\ref{thm:apu}
gives $3b_3(R)\ge4b_2(R)$, so $3m\ge3g\ge2r$.

It remains to consider a vertex-cover edge $ab$. The four-set
$L=V\setminus\{a,b\}$ is independent in $R$. Put
\[
A=N_R(a)\cap L,\quad B=N_R(b)\cap L,\quad p=|A|,\quad q=|B|,
\]
and interchange $a,b$ so that $0\le q\le p\le4$. All edges are $ab$ or
join a centre to a marked leaf, so $r=1+p+q$. Since $ab$ and $L$ are
feasible, augmentation gives a feasible triple $abx$ for some $x\in L$.
Its complement is feasible by downward closure. This partition is not
counted by $g$, because its other part lies in the independent set $L$.
Therefore
\begin{equation}\label{eq:two-matroid-extra}
m\ge g+1.
\end{equation}

For $q\ge1$, the six balanced partitions separating the centres are
parametrised by two-sets $X\subseteq L$, giving parts $\{a\}\cup X$ and
$\{b\}\cup(L\setminus X)$. They contribute to $g$ exactly when $X$ meets
$A$ and $L\setminus X$ meets $B$. Consequently
\[
g\ge6-\bin{4-p}{2}-\bin{4-q}{2}.
\]
Here $\bin{4-p}{2}$ counts the two-sets $X$ that miss $A$, and
$\bin{4-q}{2}$ counts those whose complement $L\setminus X$, itself a
two-set, misses $B$; the bound is a union bound over these two failure modes.
If $p=q=1$ and $A\ne B$, direct counting gives $g=2$. The exceptional
case $p=q=1$, $A=B$ is a triangle with isolated vertices and is treated
below. All other cases with $q\ge1$ are exhausted by the following table.
The last column, together with \eqref{eq:two-matroid-extra}, proves the
desired inequality in each case.
\begin{center}
\begin{tabular}{rrrr}
\toprule
$p$ & $q$ & Lower bound $g_0$ for $g$ & $3(g_0+1)-2(1+p+q)$\\
\midrule
1 & 1 & 2 & 3\\
2 & 1 & 2 & 1\\
3 & 1 & 3 & 2\\
4 & 1 & 3 & 0\\
2 & 2 & 4 & 5\\
3 & 2 & 5 & 6\\
4 & 2 & 5 & 4\\
3 & 3 & 6 & 7\\
4 & 3 & 6 & 5\\
4 & 4 & 6 & 3\\
\bottomrule
\end{tabular}
\end{center}

If $q=0$, the graph $R$ is a star with centre $a$, nonempty leaf set $U$,
and possible isolated vertices; here $|U|=r$. On $W=V\setminus\{a\}$,
define a graph $J$ by $ux\in E(J)$ exactly when $\{a,u,x\}$ is feasible.
Each $u\in U$ has positive degree in $J$, by augmenting $au$ towards its
feasible complementary four-set. Moreover, a degree-one vertex $u\in U$
can have its unique neighbour only outside $U$. Indeed, for each $v\in U$,
augmenting $au$ towards $V\setminus\{a,v\}$ supplies a neighbour of $u$
other than $v$.

Distribute one unit of charge from each $u\in U$ equally among its incident
edges in $J$. An edge within $U$ receives at most one, since both endpoints
have degree at least two; an edge with exactly one endpoint in $U$ also
receives at most one. Thus at least $|U|$ edges of $J$ meet $U$. Each such
edge $ux$, with $u\in U$, gives a feasible balanced partition: the triple
$aux$ is feasible, and its complement lies in the feasible four-set
$V\setminus\{a,u\}$. Distinct edges give distinct unordered partitions,
identified by their part containing $a$. Therefore $m\ge|U|=r$.

Finally, suppose $R$ is a triangle on $a,b,c$ with three isolated vertices
$L$. For each triangle edge, say $ab$, augment $ab$ towards the feasible
four-set $V\setminus ac=\{b\}\cup L$. This supplies $x\in L$ with $abx$
feasible. Its complement is a subset of the feasible four-set
$V\setminus ab$. The three triangle edges yield three distinct unordered
partitions: the part with two triangle vertices identifies the edge.
Hence $m\ge3=r$. This completes all cases.
\end{proof}

\begin{proof}[Proof of Theorem~\ref{thm:two-matroid-small}]
Deletion and contraction by a feasible fixed set preserve
\eqref{eq:two-augmentation}: for contraction by $P$, apply the original
axiom to $P\cup S$ and $P\cup T$, whose cardinality difference is unchanged.
Every positive-probability coordinate restriction therefore remains a
2-matroid, and its rank cannot increase.

For rank at most four, an even conditional ground set of size $2s\ge8$
has $F_{s-1}=\varnothing$, because the complementary set would have size
$s+1\ge5$. Its antipodal-pairs inequality is automatic. The two-point
observation and Lemmas~\ref{lem:two-matroid-four}--\ref{lem:two-matroid-six}
cover the remaining even sizes. If the original ground set instead has
size at most seven, every even conditional ground set already has size at
most six, and the same lemmas apply regardless of rank.

For positive weights, the law proportional to $w(S)\ind_{\{S\in\mathcal I\}}$
has antipodal products equal to the common factor $w(V)/Z^2$ on feasible
complementary pairs. The preceding inequalities therefore give CAPP for
this law, including after every coordinate restriction. A nonempty
downset has interval rank support. The Kahn--Neiman criterion \cite{KN}
gives the asserted weighted ultra log-concavity. Nonnegative weights follow
by continuity, with the same ambient normalisation $\bin nk$.
\end{proof}

\begin{corollary}[Initial coefficients in arbitrary rank]
\label{cor:two-matroid-initial}
For any 2-matroid on $n$ ground elements and any nonnegative element
weights, its coefficients satisfy
\[
k(n-k)A_k^2\ge(k+1)(n-k+1)A_{k-1}A_{k+1}
\qquad(1\le k\le\min(3,n-1)).
\]
Thus its first $\min(5,n+1)$ binomially normalised coefficients are
log-concave, using the original ambient order $n$.
\end{corollary}
\begin{proof}
The rank-four truncation $\{S\in\mathcal I:|S|\le4\}$ is a 2-matroid on
the same ground set. An augmentation between its feasible $S,T$ has size
$|S|+1\le|T|-1\le3$, so remains in the truncation. Its coefficients through
rank four agree with the original ones. Apply
Theorem~\ref{thm:two-matroid-small} to this truncation.
\end{proof}

The six-element lemma has no rank hypothesis. The rank bound enters only
in the reduction to even conditional ground sets of size at most six.
The next unresolved step in this approach is the eight-element
antipodal-pairs inequality, which compares feasible triples and their
complementary five-sets with feasible balanced four-set partitions.

An exact obstruction prevents simply replacing the auxiliary graph by an
arbitrary three-uniform hypergraph. Partition an eight-element set as
$V=E\mathbin{\dot\cup}\{v\}\mathbin{\dot\cup}A$, with $|E|=3$ and
$|A|=4$, and take witnesses
\[
\mathcal H=\{E\}\cup\{\{v\}\cup P:P\in\binom A2\}.
\]
Every witness has a disjoint witness. Nevertheless the doubly dependent
counts are $b_3=7$ and $b_4=8$: a successful balanced partition consists
of $E\cup\{a\}$ and $\{v\}\cup(A\setminus\{a\})$ for one of the four
$a\in A$. Thus $4b_4=32<35=5b_3$. This is an obstruction to that proof
step, not a counterexample to the $2$-matroid conjecture. Indeed, its
natural complementary downset has feasible sets $\{v\}$ and $E$, but no
feasible pair $\{v,e\}$ with $e\in E$, violating
\eqref{eq:two-augmentation}. Additional augmentation structure would be
needed at the next stage.


\section{Verification and research status}

The proof is finite and symbolic. It uses exactly two imported mathematical results: Katona's intersection-shadow theorem and the Kahn--Neiman implication from the conditional antipodal-pairs property to ultra log-concavity. Lemma~\ref{lem:ordinary}, Lemma~\ref{lem:marked}, the boundary identity, conditional closure and the weight cancellation are proved above.

The accompanying verification script uses integer arithmetic. It checks the normalised two-colouring inequalities, the boundary identity and the rank inequalities on finite systems, and independently enumerates the marked-downset lemma on small ground sets. It also checks the witness-size-three counterexamples. It covers every singleton-and-pair family on at most six labelled elements, with three weight vectors each, and the marked-downset lemma for every downset on at most five elements. All checks are explicit conditionals, so they remain active under \texttt{python -O}. Finite checks cannot establish the universal theorem; they test the implementation and the vulnerable normalisations and boundary cases. The exact scope and results are in \texttt{verification\_results.json}.

A separate complete six-element enumeration corroborates the $2$-matroid lemma. Its reduction, readable C++ source and exact results are supplied in the supplementary directory. It checked $2{,}653{,}437$ four-set/triple-family combinations and $1{,}723{,}681$ valid maximal-pair completions, finding no negative antipodal margin. Both the main graph theorem and the $2$-matroid extension have analytic proofs; the enumeration is not a proof dependency.

The bounded literature search was performed on 2 November 2026 and repeated on 4 November 2026, including works citing Xie and Xu; no proof of Conjecture~4.1 was found. It inspected the October 2026 version of the originating conjecture and searched equivalent formulations involving noncover complexes, Alexander duals, codimension-two complexes, monotone Boolean formulas, weighted rank sequences and antipodal pairs. No prior proof of the stated weighted theorem or normalised doubly-dependent theorem was found within that search. Negative search evidence does not establish historical priority. The accompanying report records the search boundary and the sources inspected.

\paragraph{Status.} This is an AI-assisted research paper. It has had several internal proof audits, including an independently reconstructed proof and adversarial review, but no external peer review and no proof-assistant formalisation. The natural next step, the eight-element antipodal inequality needed for $2$-matroids of rank five, meets the explicit obstruction described in Section~\ref{sec:two-matroid}; it was not attempted here.

\paragraph{Archive.} This version (0.1.0-candidate) is archived with DOI \href{https://doi.org/10.5281/zenodo.23935650}{10.5281/zenodo.23935650} (all versions: \href{https://doi.org/10.5281/zenodo.23935649}{10.5281/zenodo.23935649}). The repository is \url{https://github.com/ipitchford/dependence-ultra-log-concavity}.

\begingroup\small
\begin{thebibliography}{99}
\bibitem{XX}
Xie, Y.-T., \& Xu, S.-J. (2026). Ultra log-concavity and real-rootedness of dependence polynomials. \emph{The Electronic Journal of Combinatorics, 33}(4), P4.9. \url{https://doi.org/10.37236/13310}. Inspected version: \url{https://arxiv.org/html/2408.09152v2}, 10 October 2026. The preprint first appeared in 2024.

\bibitem{H02}
Horrocks, D. G. C. (2002). The numbers of dependent $k$-sets in a graph are log concave. \emph{Journal of Combinatorial Theory, Series B, 84}(1), 180--185. \url{https://doi.org/10.1006/jctb.2001.2077}. Abstract and bibliographic record inspected; full text not obtained; no technical step here depends on its proof.

\bibitem{H00}
Horrocks, D. G. C. (2000). Doubly dependent sets in graphs. \emph{Australasian Journal of Combinatorics, 22}, 105--116. \url{https://ajc.maths.uq.edu.au/pdf/22/ocr-ajc-v22-p105.pdf}.

\bibitem{Katona}
Katona, G. O. H. (1964). Intersection theorems for systems of finite sets. \emph{Acta Mathematica Academiae Scientiarum Hungaricae, 15}, 329--337. Repository record: \url{https://real.mtak.hu/21124/}. The theorem used here was inspected in Frankl's primary exposition below.

\bibitem{Frankl}
Frankl, P. (2016). Linear independence, a unifying approach to shadow theorems. \emph{European Journal of Combinatorics, 57}, 36--39. \url{https://doi.org/10.1016/j.ejc.2016.04.001}. Full text: \url{https://www.renyi.hu/~pfrankl/2016-3.pdf}.

\bibitem{ALOV}
Anari, N., Liu, K., Oveis Gharan, S., \& Vinzant, C. (2024). Log-concave polynomials III: Mason's ultra-log-concavity conjecture for independent sets of matroids. \emph{Proceedings of the American Mathematical Society}. \url{https://doi.org/10.1090/proc/16724}. Preprint: arXiv:1811.01600.

\bibitem{BH}
Br\"and\'en, P., \& Huh, J. (2020). Lorentzian polynomials. \emph{Annals of Mathematics, 192}(3), 821--891. \url{https://doi.org/10.4007/annals.2020.192.3.4}.

\bibitem{KK}
Kim, J., \& Kim, M. (2021). Domination numbers and noncover complexes of hypergraphs. arXiv:2101.01850. \url{https://arxiv.org/abs/2101.01850}.

\bibitem{Liggett}
Liggett, T. M. (1997). Ultra logconcave sequences and negative dependence. \emph{Journal of Combinatorial Theory, Series A, 79}(2), 315--325. \url{https://doi.org/10.1006/jcta.1997.2790}.

\bibitem{KN}
Kahn, J., \& Neiman, M. (2011). A strong log-concavity property for measures on Boolean algebras. \emph{Journal of Combinatorial Theory, Series A, 118}(6), 1749--1760. \url{https://doi.org/10.1016/j.jcta.2011.02.007}. Full preprint inspected: \url{https://arxiv.org/pdf/0907.0243}.
\end{thebibliography}
\endgroup
\end{document}
